\chapter{Ventilation-Perfusion (V/Q) Scan}

\section*{What Is This Test?}
A V/Q scan has two parts: ventilation (inhaling radioactive gas) and perfusion (injected radioactive material). It compares air flow and blood flow in the lungs. A mismatch (ventilation normal, perfusion absent) suggests pulmonary embolism.

\section*{Alternative Names}
V/Q scan, lung perfusion scan, lung ventilation scan, ventilation-perfusion scintigraphy, pulmonary V/Q


\section*{Principle}
The V/Q scan uses two tracers: ventilation imaging uses inhaled radioactive gas (Xe-133 or Kr-81m) or aerosol (Tc-99m DTPA); perfusion imaging uses injected Tc-99m macroaggregated albumin (MAA) that lodges in pulmonary capillaries. A gamma camera images the distribution. Perfusion defects with preserved ventilation indicate pulmonary embolism (mismatch).

\section*{History}
The V/Q scan was developed in the 1960s using radioactive gases and macroaggregated albumin. It became the standard noninvasive test for pulmonary embolism diagnosis until the widespread adoption of CT angiography.

\section*{Why This Test Is Ordered}
Your doctor orders a V/Q scan when they suspect a pulmonary embolism (blood clot in your lungs) and a CT scan cannot be performed --- for example, if you have a contrast allergy or kidney disease. It answers the question: is there a mismatch between air flow and blood flow in your lungs that would indicate a blood clot?

\section*{Is This a Primary or Secondary Test}
A secondary test. CT pulmonary angiography is the preferred first test for suspected pulmonary embolism. V/Q scanning is used when CT is not possible.

\section*{How This Test Is Performed}
You inhale a small amount of radioactive gas or aerosol through a mask (ventilation part), and receive a small injection of radioactive material into your vein (perfusion part). A special camera then takes images of your lungs. The entire test takes thirty to sixty minutes.

\section*{What Preparation Is Needed}
No special preparation is needed. Tell your doctor if you are pregnant or breastfeeding.

\section*{Contraindications and Precautions}
V/Q scanning uses radioactive tracers and is generally not performed during pregnancy. Inform the team if you are pregnant or breastfeeding. The ventilation portion may be difficult for patients with severe respiratory distress.

\section*{Risks}
V/Q scanning involves a small dose of radioactive material, with radiation exposure comparable to a chest CT scan. Rare allergic reactions to the tracer may occur. The inhaled ventilation tracer may cause mild throat irritation. Inform the team if you are pregnant or breastfeeding.

\section*{What Happens After the Test}
You can resume normal activities immediately. Results are typically available within one to three days.

\section*{Understanding Your Results}

\section*{Normal Results}
A normal or low-probability result effectively rules out pulmonary embolism.

\section*{Follow-up Tests}
If the V/Q scan is positive for pulmonary embolism, your doctor will start anticoagulation therapy. If the result is indeterminate, a CT pulmonary angiography may be performed. If the V/Q scan is negative, other causes of your symptoms will be investigated.
\clearpage
